% TOP_PROBLEM: {"id": 30004562, "title": "Bilinear Forms behind Resistor and Ising Integrable Equations", "category_id": 16, "category": {"id": 16, "name": "physics", "display_name": "Mathematical Physics", "description": "Problems at the intersection of mathematics and physics.", "slug": "physics", "order_index": 16, "created_at": "2026-07-31T15:26:25.671Z"}, "status": "open", "statusQualification": "The relevant subvarieties, diagram maps and integrable equations are source-defined, not arbitrary bilinear forms on arbitrary vector spaces. The question asks for the mechanism behind the reported examples, not simply verification of their names."}
% FIRST_POSTED: 2026-10-11
% ENTRY_KIND: statement_only
% SUBMISSION: {"kind": "statement_only", "category": "top_problems", "problem_id": "30004562", "proposal_type": "established", "source": "Mini-Workshop: Dimers, Ising and Spanning Trees beyond the Critical Isoradial Case (2020), Oberwolfach Reports, DOI 10.4171/owr/2020/35, pp. 1759--1789. Primary report accessed 2026-10-11; relevant PDF pages 22. https://doi.org/10.4171/owr/2020/35; Affolter, discrete differential geometry and cluster algebras (catalogue-listed context; not independently consulted in this preparation). https://arxiv.org/abs/2305.02212", "status_evidence": "Canonical source status open retained. Supplied literature triage dated 2026-08-22: The symmetric-bilinear-to-resistor/BKP and antisymmetric-to-Ising/CKP pattern is documented in examples, but a general conceptual explanation or theorem identifying the mechanism is still an open question. Primary-report access checked 2026-10-11; this historical report and inherited triage do not certify current resolution of every component.", "status_checked_on": "2026-10-11", "duplicate_check": "Compared canonical ID, title and normalized target with the complete local UnsolvedMath registry; normalized exact-text matches: none; title matches: none. Related or overlapping canonical records: none identified in this batch. Primary source context and the existing OWR statement audit used for scope and duplicate review; an exhaustive mathematical-equivalence/literature search was not performed.", "attribution": "Original proposer not named in the supplied canonical record; attribution remains to the identified Oberwolfach contribution/report, not its editors as invented proposers. Explanatory statement preparation and source/semantic audit assisted by OpenAI GPT-6 Codex; no human mathematical authorship or independent proof verification is claimed."}
% SOURCE_STATEMENT: "Why do symmetric bilinear forms lead to the resistor subvariety and the BKP equation, whereas antisymmetric bilinear forms lead to the Ising subvariety and the CKP equation?"
% Preparation and final self-audit: CONTRIBUTING.md prompts C and D; source/audit artifacts retained in the local physics_statements_2026_10_11 batch.
% Imported from: physics_statements.tex; UnsolvedMath ID 30004562
% !TeX program = lualatex
\documentclass[11pt]{article}
\usepackage[margin=25mm]{geometry}
\usepackage{fontspec}
\setmainfont{Latin Modern Roman}
\usepackage{amsmath,amssymb,mathtools}
\usepackage{unicode-math}
\setmathfont{Latin Modern Math}
\usepackage{xurl,hyperref}
\hypersetup{colorlinks=true,urlcolor=blue}
\setcounter{secnumdepth}{0}
\setlength{\parindent}{0pt}
\setlength{\parskip}{5pt}
\setlength{\emergencystretch}{3em}
\newcommand{\sref}[2]{\hyperlink{source:#1:#2}{[S#2]}}
\newcommand{\cataloguescope}[1]{\paragraph{Recorded target.}#1}
\newcommand{\R}{\mathbb R}
\newcommand{\C}{\mathbb C}
\newcommand{\N}{\mathbb N}
\newcommand{\Z}{\mathbb Z}
\newcommand{\Q}{\mathbb Q}
\newcommand{\F}{\mathbb F}
\newcommand{\eps}{\varepsilon}
\begin{document}
% UnsolvedMath ID 30004562; source catalogue code OWR-2654831-005
\setcounter{section}{30004561}
\section{Bilinear Forms behind Resistor and Ising Integrable Equations}\label{30004562}
{\small UnsolvedMath ID 30004562 · Mathematical Physics}
\subsection{Definitions and mathematical statement}

\paragraph{Shared notation.}$\mathbb N=\{1,2,\ldots\}$; the standard number fields have their usual meanings. All additional conventions are stated locally.


\paragraph{Definitions and conventions.}
A bilinear form $B$ is symmetric when $B(u,v)=B(v,u)$ and antisymmetric when $B(u,v)=-B(v,u)$. The report\textquotesingle s resistor and Ising subvarieties are constrained spaces of variables for triple-crossing-diagram maps in projective geometry. BKP and CKP are the named reductions or equations in the integrable-system hierarchy. The question concerns why the symmetry type of the geometric form selects one of those network and integrability structures.

\paragraph{Target retained from the canonical source catalogue.}
\paragraph{Catalogue wording retained for traceability.}
Why do symmetric bilinear forms lead to the resistor subvariety and the BKP equation, whereas antisymmetric bilinear forms lead to the Ising subvariety and the CKP equation?

\paragraph{Source scope and qualification.}
The relevant subvarieties, diagram maps and integrable equations are source-defined, not arbitrary bilinear forms on arbitrary vector spaces. The question asks for the mechanism behind the reported examples, not simply verification of their names.
\subsection{Short English statement}
Explain conceptually why symmetric and antisymmetric projective bilinear forms produce the indicated resistor/Ising and BKP/CKP pairings.
\subsection{Sources}
\begin{itemize}
\item[\hypertarget{source:30004562:1}{S1}] ULAM.AI and the UnsolvedMath Contributors, UnsolvedMath, record 30004562 (OWR-2654831-005). Catalogue curation under CC BY 4.0; definitions and English exposition added locally.
\url{https://huggingface.co/datasets/ulamai/UnsolvedMath}.
\item[\hypertarget{source:30004562:2}{S2}] Mini-Workshop: Dimers, Ising and Spanning Trees beyond the Critical Isoradial Case (2020), Oberwolfach Reports, DOI 10.4171/owr/2020/35, pp. 1759--1789. Primary report accessed 2026-10-11; relevant PDF pages 22.
\url{https://doi.org/10.4171/owr/2020/35}.
\item[\hypertarget{source:30004562:3}{S3}] Affolter, discrete differential geometry and cluster algebras (catalogue-listed context; not independently consulted in this preparation).
\url{https://arxiv.org/abs/2305.02212}.
\end{itemize}
\label{end:30004562}
\end{document}
